\section{Problem and deployment constraints}
\label{sec:problem}
An image enhancer maps a low-light observation $x_t$ to an enhanced frame
$z_t=f_\theta(x_t)$. A video system must also preserve legitimate scene changes,
avoid noise-driven fluctuations, and fit its computation and memory budget.
Low parameter count alone does not establish any of these properties.
MobileIE~\cite{mobileie} provides a useful compact starting point. Our
12-channel low-light deployment instance has 4,075 parameters; \emph{4K parameters}
does not mean a video resolution of $4\mathrm{K}$.

We distinguish a stateless predictor from a stateful one:
\[
 z_t=f_\theta(x_t),\qquad
 (z_t,h_t)=F_\theta(x_t,h_{t-1}).
\]
A training objective may use neighbouring frames without introducing
$h_{t-1}$ at deployment. Conversely, a computation whose output changes for
two histories with the same current frame cannot generally be represented
exactly by a function of that current frame alone. Spatial convolution/BN
folding does not eliminate this information dependency.

Our evaluation therefore reports quality, temporal error, parameters,
convolution operations, online history, and measured pipeline latency as
separate quantities. A small parameter change can still introduce
full-resolution arithmetic or memory traffic. Inference motion estimation,
sampling, padding, and buffering belong to the deployed pipeline even when
they have no learned parameters. Server CPU measurements are not phone
measurements.

\subsection{The original nonlinear pathway}
The retained MobileIE backbone consists of a convolutional head, a second
feature block, channel/spatial gates, and an RGB readout~\cite{mobileie}.
Its feature self-transform (FST) has the form
\[
 q(u)=(w_1u)(w_2u)+b=a u^2+b,
\]
where $w_1,w_2$ are shared scalars and $b$ is channel-wise. Channel attention
uses global average pooling followed by a $12\!\times\!12$ pointwise map and
a sigmoid. The spatial gate depends on the channel-weighted feature maximum.
The learned enhancement mapping, global statistics, and local nonlinear
response can consequently interact.

For a scalar perturbation $\epsilon$, the identity
\[
 q(u+\epsilon)-q(u)=2au\epsilon+a\epsilon^2
\]
motivates examining noise sensitivity at the root operator. If the
perturbation has conditional mean zero, its conditional mean shift is
$a\,\mathbb E[\epsilon^2]$. This elementary identity is not a new
contribution. It neither proves that FST is the sole cause of video flicker
nor establishes an unbiased correction for features after nonlinear
convolutions, ISP processing, clipping, and quantization.

\subsection{Known mechanisms and the remaining research question}
Training-time reparameterization, channel multiplication, temporal
regularization, and classical motion-compensated fusion are established
tools. FastDVDnet~\cite{fastdvdnet}, for example, learns multiframe denoising
without an explicit flow estimator. TLC~\cite{tlc} addresses a different issue:
the spatial-statistic mismatch between patch training and whole-image
inference. Learnable safe rational activations also have clear
precedents~\cite{pau}.

The unresolved question is whether an independently justified video
adaptation can improve the quality--stability--cost trade-off over strong
matched controls in this extremely compact setting. Neither assembling these
known components nor obtaining a nonzero historical effect answers that
question. The current candidate remains an experiment until its controlled
results, external validity, and deployment costs support a defensible claim.
